\documentclass{alltagsforschung}

\journalvolume{7}
\journalissue{1}
\journalyear{2026}
\articledate{04.11.2026}
\articlepages{8--17}

\title{Zur empirischen Bedeutung des Wartens vor dem Kühlschrank}

\subtitle{Eine alltagswissenschaftliche Untersuchung spontaner Öffnungsvorgänge}

\author{Prof. Dr. Hildegard Wankelmuth}

\shorttitlecommand{Warten vor dem Kühlschrank}

\articlekeywords{
  Alltag; Kühlschrank; Entscheidungsverhalten; Beobachtungsstudie
}

\abstracttext{
Untersucht wird das wiederholte Öffnen eines Kühlschranks in Situationen,
in denen keine konkrete Veränderung seines Inhalts zu erwarten ist.
Die Untersuchung zeigt, dass die Handlung trotz fehlender neuer Information
regelmäßig wiederholt wird.
}

\begin{document}

\section{Ausgangslage}

Das Öffnen des Kühlschranks stellt eine alltägliche Handlung dar, deren
Notwendigkeit häufig unmittelbar nachvollziehbar ist.

Weniger eindeutig ist die Situation, in der der Kühlschrank kurz nach dem
Schließen erneut geöffnet wird, obwohl zwischenzeitlich keine relevante
Veränderung des Inhalts stattgefunden haben kann.

\section{Fragestellung}

Die Studie untersucht, unter welchen Bedingungen ein solcher Vorgang
beobachtet werden kann und welche Erwartungen mit ihm verbunden sind.

\section{Beobachtung}

In mehreren Beobachtungssituationen wurde festgestellt, dass der Vorgang
auch ohne konkrete Erwartung eines neuen Lebensmittels wiederholt wird.

\section{Ergebnis}

Die Wiederholung scheint nicht ausschließlich durch die Erwartung eines
veränderten Kühlschrankinhalts erklärt werden zu können.

\section{Diskussion}

Eine abschließende Erklärung des Phänomens konnte im Rahmen der vorliegenden
Untersuchung nicht gefunden werden.

Dies ist insofern bemerkenswert, als der Kühlschrank selbst keine
Veränderung erkennen ließ.

\end{document}
